\documentclass[tikz,border=10pt]{standalone}
\usepackage{ctex}
\usepackage{amsmath}
\usetikzlibrary{calc, arrows.meta}

\begin{document}
% geo-p9-05-2: 二面角法向量法 + 同指 vs 异指 判正负
\begin{tikzpicture}[
  x={(1cm,0cm)},
  y={(0.40cm,0.32cm)},
  z={(0cm,1cm)},
  thick,
  >=Stealth,
  line cap=round, line join=round,
]

%% ============================================================
%% 左栏：二面角为锐角（法向量"同侧"，cos > 0）
%% ============================================================
\begin{scope}[xshift=0cm]
  \node[font=\bfseries\small] at (1.5, -0.5, 3.8) {情形 A：锐角二面角};
  \node[font=\small] at (1.5, -0.5, 3.3) {（法向量指向同侧，$\cos\psi > 0$）};

  % 棱
  \coordinate (E1) at (1.5, 0.0, 0);
  \coordinate (E2) at (1.5, 3.0, 0);
  \draw[black, very thick] (E1) -- (E2) node[right, font=\small] {棱 $l$};

  % 半平面 α（浅蓝，水平）
  \fill[blue!15, opacity=0.85]
    (E1) -- (E2) -- (-0.5, 3.0, 0) -- (-0.5, 0.0, 0) -- cycle;
  \node[blue!70, font=\small] at (-0.1, 2.0, 0.2) {$\alpha$};

  % 半平面 β（浅绿，小角度向上倾斜）
  \fill[green!15, opacity=0.85]
    (E1) -- (E2) -- (3.5, 3.0, 1.5) -- (3.5, 0.0, 1.5) -- cycle;
  \node[green!60!black, font=\small] at (3.0, 2.0, 1.3) {$\beta$};

  % 棱上基准点 O
  \coordinate (O1) at (1.5, 1.5, 0);
  \fill[black] (O1) circle [radius=1.5pt];
  \node[left, font=\small] at (O1) {$O$};

  % α 的法向量 n1（竖直向上，指向 α 面上方）
  \draw[blue!80, very thick, ->] (O1) -- ++(0, 0, 2.2)
    node[above, font=\small, blue!80] {$\vec{n_1}$};

  % β 的法向量 n2（稍向左偏，也指向上方——两者"同侧"）
  \draw[green!70!black, very thick, ->] (O1) -- ++(-0.4, 0, 2.2)
    node[above left, font=\small, green!70!black] {$\vec{n_2}$};

  % 两法向量夹角弧（与二面角一致，锐角）
  \draw[red!60, thick]
    ($(O1) + (0, 0, 0.9)$) arc[start angle=90, end angle=100, radius=0.9cm];

  % 二面角平面角弧（小角，锐）
  \draw[red!70!black, thick]
    ($(O1) + (-0.8, 0, 0)$) arc[start angle=180, end angle=148, radius=0.8cm]
    node[midway, left=2pt, font=\small, red!80!black] {$\psi$（锐）};

  % 结论
  \node[font=\small, align=center, text width=4cm] at (1.5, -0.5, -1.2)
    {$\vec{n_1},\vec{n_2}$ 指向同侧\\[2pt]$\Rightarrow\cos\psi=\dfrac{\vec{n_1}\cdot\vec{n_2}}{|\vec{n_1}||\vec{n_2}|}>0$\\[2pt]取正值};
\end{scope}

%% ============================================================
%% 右栏：二面角为钝角（法向量"异侧"，cos < 0）
%% ============================================================
\begin{scope}[xshift=7.5cm]
  \node[font=\bfseries\small] at (1.5, -0.5, 3.8) {情形 B：钝角二面角};
  \node[font=\small] at (1.5, -0.5, 3.3) {（法向量指向异侧，$\cos\psi < 0$）};

  % 棱
  \coordinate (E1) at (1.5, 0.0, 0);
  \coordinate (E2) at (1.5, 3.0, 0);
  \draw[black, very thick] (E1) -- (E2) node[right, font=\small] {棱 $l$};

  % 半平面 α（浅蓝，水平）
  \fill[blue!15, opacity=0.85]
    (E1) -- (E2) -- (-1.2, 3.0, 0) -- (-1.2, 0.0, 0) -- cycle;
  \node[blue!70, font=\small] at (-0.5, 2.0, 0.2) {$\alpha$};

  % 半平面 β（浅绿，大角度向上倾斜——钝角）
  \fill[green!15, opacity=0.85]
    (E1) -- (E2) -- (4.0, 3.0, 3.2) -- (4.0, 0.0, 3.2) -- cycle;
  \node[green!60!black, font=\small] at (3.5, 2.0, 3.0) {$\beta$};

  % 棱上基准点 O
  \coordinate (O2) at (1.5, 1.5, 0);
  \fill[black] (O2) circle [radius=1.5pt];
  \node[left, font=\small] at (O2) {$O$};

  % 二面角弧（钝角）
  \draw[red!70!black, thick]
    ($(O2) + (-0.9, 0, 0)$) arc[start angle=180, end angle=115, radius=0.9cm]
    node[midway, left=2pt, font=\small, red!80!black] {$\psi$（钝）};

  % α 法向量 n1（向上，指向 α 上方）
  \draw[blue!80, very thick, ->] (O2) -- ++(0, 0, 2.2)
    node[above, font=\small, blue!80] {$\vec{n_1}$};

  % β 法向量 n2（垂直于 β 面，指向 β 面一侧，即向右偏下——与 n1 异侧）
  % β 面方向向量 (1, 0, 1.3)，法向量 = (-1.3, 0, 1)（旋转90度）——指向 α 那侧
  \draw[green!70!black, very thick, ->] (O2) -- ++(-1.4, 0, 1.1)
    node[above left, font=\small, green!70!black] {$\vec{n_2}$};

  % 两法向量夹角弧（大于90度，钝角）
  \draw[orange!80!black, dashed, thick]
    ($(O2) + (0, 0, 1.0)$) arc[start angle=90, end angle=141, radius=1.0cm]
    node[midway, above=2pt, font=\tiny, orange!80!black] {钝角};

  % 结论
  \node[font=\small, align=center, text width=4.5cm] at (1.5, -0.5, -1.2)
    {$\vec{n_1},\vec{n_2}$ 指向异侧\\[2pt]$\Rightarrow\cos(\text{法向量夹角})<0$\\[2pt]$\Rightarrow$二面角 $=\pi -$ 法向量夹角\\[2pt]取对应符号};

  % 判断规则框
  \node[font=\small, align=left, text width=5.5cm,
        draw=gray, rounded corners=3pt, fill=yellow!10,
        inner sep=4pt] at (1.5, -0.5, -3.2) {
    \textbf{判断规则}\\
    法向量 $\vec{n_1},\vec{n_2}$ 指同侧 $\Rightarrow$ 二面角锐，$\cos>0$\\
    法向量 $\vec{n_1},\vec{n_2}$ 指异侧 $\Rightarrow$ 二面角钝，$\cos<0$\\
    不确定时：看图判断后取对应符号
  };
\end{scope}

\end{tikzpicture}
\end{document}
